\documentclass[10pt]{article}
\usepackage[margin=20mm]{geometry}
\usepackage{amsmath,amssymb,mathtools,fontspec,unicode-math,xurl,hyperref}
\IfFontExistsTF{DejaVu Serif}{\setmainfont{DejaVu Serif}}{\setmainfont{Latin Modern Roman}}
\setmathfont{latinmodern-math.otf}
\setlength{\parindent}{0pt}\setlength{\parskip}{6pt}\sloppy
\title{Initial scientific review of
Characterizing Gibbs states for area-tilted Brownian lines}
\author{Math2025 review record M25-0533}
\date{29 October 2026}
\begin{document}
\maketitle
Abstract: This review identifies one false supporting estimate in the multiline classification proof and one peripheral overstatement of uniformity. The heavy-path envelope of Lemma 6.8 has an incorrect boundary-layer width; an admissible parameter family makes its asserted high-probability event have probability tending to zero. The all-positive-tilt version of (2.3) also fails under its own scaling relation. These conclusions concern the stated estimates and their uses. They do not disprove the final classification theorems. Possible repairs, counterevidence and the limits of this initial review are stated separately.\par
\section*{Reviewed source, framework and assessment}
The reviewed paper is by Mriganka Basu Roy Chowdhury, Pietro Caputo and Shirshendu Ganguly, The Annals of Probability 53(6) (2025), 2196–2255, DOI 10.1214/25-AOP1760 (https://doi.org/10.1214/25-AOP1760) [BCG]. The frozen source is the 60-page institutional PDF designated as the version of record by the acquisition manifest. Its SHA256 begins `df44a62d47f5de6e`; the complete digest, exact bytes and acquisition origin are retained in the accompanying provenance. The institutional record labels the file an editorial version. A separate byte-equivalence claim about every publisher-hosted copy is not made. Physical PDF page numbers are used throughout for source locations.\par
The primary scientific executor read all 60 source pages, including the definitions, proof arguments and references. The source image checks include all pages supporting the findings below. Reference reading was targeted, with the exact acquired versions and page subsets recorded; it was not a complete review of the cited literature. The assessment retains the initial review scope, with the interval-maximum candidate finding withdrawn after independent verification of the source formula and its length factor. Scientific disposition remains pending and no independent acceptance or human review is represented by this report.\par
\subsection*{Definitions and hypotheses retained}
Equation reference: eq:pbr. The one-line measure \(\mathcal L^{a;x,y}_{\ell,r}\) is the Brownian bridge from \(x>0\) to \(y>0\) on \([\ell,r]\), weighted by \(\exp(-a\int_\ell^r X(s)\,ds)\) and conditioned to remain positive (Definition 1.1, PDF p. 4). Thus, for strength \(2a\), the parabolic barrier representation is \[
 X(t)=Y(t)+at^2,\qquad
 Y\ \text{a standard Brownian bridge above }-at^2.\] This normalization follows directly from Lemma 2.4, PDF p. 13; using strength \(a\) instead would give the parabola \(at^2/2\). The Brownian variance convention is the one in the paper.\par
In the infinite ensemble, \(\lambda>1\) is fixed, the area strengths are \(2\lambda^{i-1}\), the curves are strictly ordered and positive, and the Brownian–Gibbs property resamples only finitely many curves on a bounded interval. Definition 1.7 additionally requires asymptotic pinning to zero. This assumption cannot be dropped: an upward translation preserves the Gibbs specification but changes the floor selected at infinity (PDF p. 8).\par
Theorems 1.3 and 1.8 classify extremal states by \((L,R)\in[-\infty,\infty)^2\) with \(L+R<0\); finite slopes, one slope equal to \(-\infty\), and both slopes equal to \(-\infty\) are distinct cases. Positive infinity is excluded. The parabolic coefficient is fixed by the tilt normalization. Time translation changes the slopes by \((L,R)\mapsto(L+2h,R-2h)\) at strength 2. Theorem 1.9 concerns lower lines and includes the time-scale term \(\lambda^{2k/3}\) inside the interval-maximum expectation. None of the counterexamples below violates these distinctions.\par
\subsection*{Proof dependencies}
For the multiline argument, the decisive chain runs from the heavy and light path estimates (Lemmas 6.1–6.3) to Proposition 5.5, then to confinement in Lemma 5.6, top-line fluctuation control in Lemma 5.7 and the reverse coupling used for uniqueness. Lemma 6.8 is explicitly applied in the heavy-path induction (PDF pp. 56–58). Equation (5.9) is also invoked in Lemma 5.7 (PDF p. 40), but the source formula includes the factor \(S\). That factor is precisely the interval-length dependence expected from the union-bound step; the initial report omitted it when transcribing the claim. No defect is asserted for (5.9). The remaining false supporting estimate does not by itself establish a false final classification theorem.\par
\section*{Problem 1: the heavy-path envelope starts too early}
\subsection*{Locator, assertion and hypotheses}
Equation reference: eq:printedwidth. Lemma 6.8, PDF pp. 56–57, journal pp. 2251–2252, takes \[X\sim\mathcal L_T^{2a;H},\qquad a=\lambda^{k-1},\quad k>n_0(T),
 \quad H\leq 2T^2,\quad 0<\delta<1/20,\] where \(n_0(T)\) is the largest integer \(j\) such that \(T^{-15}\leq T\lambda^{-(j-1)/2}\) (PDF p. 47). The asserted envelope is \(H+\Delta_2\) in the two endpoint strips and \(\Delta_2\) in the central interval, with \[
 \Delta_1=T\lambda^{-k/2},\qquad
 \Delta_2=T^\delta\lambda^{-(k-n_0)/5}.\] Its probability is claimed to be at least \(1-\exp(-\Theta(\lambda^{k\delta/2}))\). The failure below occurs strictly inside the central interval, so the overlap of the piecewise definition at its endpoints is immaterial.\par
\subsection*{An admissible family and an exact lower comparison}
Equation reference: eq:parameters, eq:tangent. Fix \[
 \lambda=4,\quad \delta=\frac1{40},\quad T=2^m,\quad m\geq2,
 \quad n_0=16m+1,\quad k=16m+2,\quad H=T^2.\] The formula for \(n_0\) follows by comparing powers of 2 in its defining inequality. All the hypotheses above hold, with a positive boundary height. Put \[d=\Delta_1=T2^{-k}=\frac14T^{-15},
 \qquad a=4^{k-1}.\] Then \(ad^2=T^2/4\) and the true tangency point from \((T,H-aT^2)\) to the barrier \(-at^2\) is \[
 \xi=T-\sqrt{H/a}=T-2d.\] Consider \(t_0=T-5d/4\). It satisfies \(\xi<t_0<T-d\) and belongs to the central interval in the printed envelope.\par
Equation reference: eq:mean, eq:variance. Condition on \(Y(\xi)\) in the representation [eq:pbr] (\#eq:pbr). The bridge on \([\xi,T]\) has endpoint \(H-aT^2\) at \(T\) and \(Y(\xi)>-a\xi^2\) at its left endpoint. Lowering that endpoint to \(-a\xi^2\) and removing the lower-barrier constraint gives a stochastically smaller, ordinary Brownian bridge. Its mean is the tangent line to \(-at^2\) at \(\xi\). After adding back \(at^2\), the resulting Gaussian lower comparison for \(X(t_0)\) has \[\begin{aligned}
 \mu_m&=a(t_0-\xi)^2=\frac9{64}T^2, \\
 \sigma_m^2&=\frac{(t_0-\xi)(T-t_0)}{T-\xi}
           =\frac{15}{32}d=\frac{15}{128}T^{-15}.
 \end{aligned}\] Both comparisons are in the required direction: increasing an endpoint or imposing a lower floor increases the bridge law. The comparison therefore holds after averaging over the conditioned value \(Y(\xi)\).\par
Proposition 1.  For the family [eq:parameters] (\#eq:parameters), the probability of the envelope event in Lemma 6.8 tends to zero as \(m\to\infty\).\par
Equation reference: eq:probupper, eq:chernoff. Proof. At \(t_0\) the envelope requires \(X(t_0)\leq
T^{1/40}4^{-1/5}\). Hence, with \(\Phi\) the standard Gaussian distribution function, \[
 \mathbb P(X\leq\mathrm{Envelope})
 \leq
 \Phi\!\left(
 \frac{T^{1/40}4^{-1/5}-9T^2/64}
      {\sqrt{15T^{-15}/128}}\right)\longrightarrow0.\] For \(T\geq4\), the threshold is at most \(\mu_m/2\). The elementary Gaussian Chernoff bound consequently gives the more explicit estimate \[
 \mathbb P(X\leq\mathrm{Envelope})
 \leq \exp\!\left(-\frac{27}{1280}T^{19}\right).\] In contrast, \(\lambda^{k\delta/2}=2^{1/20}T^{2/5}\to\infty\), so the claimed lower probability tends to one for every positive constant hidden in \(\Theta\). ◻\par
This is a probabilistic counterexample proved by exact Gaussian comparison, not a simulation. Independent symbolic checks verified the tangency identity, mean, variance, threshold index and Chernoff coefficient. For example, \(T=16\) gives \(n_0=65\), \(k=66\), \(d=2^{-62}\), \(\mu_m=36\), \(\sigma_m^2=15/2^{67}\), and a printed bulk ceiling approximately \(0.8122523964\).\par
\subsection*{The actual induction also encounters the defect}
Equation reference: eq:inductionmean. The application in the proof of Lemma 6.1, PDF p. 58, uses boundary height \(H=T^2-2KT\), not an arbitrary unused endpoint of the lemma’s range. The same test point works for every fixed \(K>0\) and all sufficiently large \(T\) in [eq:parameters] (\#eq:parameters). Its tangency width is \[\sqrt{H/a}=2d\sqrt{1-2K/T}.\] The Gaussian comparison now has mean \[
 T^2\left(\sqrt{1-2K/T}-\frac58\right)^2
 \sim\frac9{64}T^2,\] and variance of order \(d\). Thus it again exceeds the claimed bulk ceiling with probability tending to one. This rules out the defense that only irrelevant values near \(2T^2\) caused the contradiction.\par
\subsection*{Counterevidence, repair and impact}
The proof itself locates tangency only up to \(\xi=T(1-O(\lambda^{-k/2}))\). That statement does not put tangency after the specific point \(T-\Delta_1\) with coefficient 1. For all \(H\leq2T^2\), the geometric boundary-layer width is bounded by \[\sqrt{H/a}\leq\sqrt{2\lambda}\,T\lambda^{-k/2}.\] A plausible repair is to replace \(\Delta_1\) by \(C_\lambda T\lambda^{-k/2}\) with \(C_\lambda>\sqrt{2\lambda}\) and then propagate that constant through the induction. This is consistent with the proof’s order-of-magnitude geometry. For fixed \(\lambda>1\), the cumulative heavy-layer widths remain \(O_\lambda(T^{-15})\), which is smaller than \(T^{-10}\) for large \(T\). Thus the downstream domain margin does not itself obstruct such a repair.\par
This observation is counterevidence against a claim that the final theorem is false. It is not a substitute proof of the repaired probability estimate: the excursion comparison, fluctuation allowance and substitution of shortened domains in the induction must still be checked with the corrected constant. The finding is Major: a false lemma is used explicitly to establish the heavy-path input to Proposition 5.5. No assertion of failure of Theorems 1.3, 1.8 or 1.9 follows.\par
\section*{Problem 2: the Ferrari–Spohn maximum bound overstates small-tilt uniformity}
\subsection*{Locator and an admissible scaling limit}
Equation reference: eq:badFS, eq:smalllower. Lemma 2.6(5), equation (2.3), PDF p. 15, journal p. 2210, claims constants \(C,c>0\) independent of all \(T\geq1\), \(t>0\), \(a>0\) such that, for \(X_a\sim\mu^{a,\mathrm{FS}}\), \[
 \mathbb P\!\left(\max_{s\in[-T,T]}X_a(s)>t\right)
 \leq CTa^{2/3}\exp(-c\sqrt a\,t^{3/2}).\] The same lemma’s item 3 gives the exact scaling, with \(X\sim\mu^\mathrm{FS}=\mu^{2,\mathrm{FS}}\), \[X_a(s)\ \stackrel{\rm law}{=}\
 (a/2)^{-1/3}X\big((a/2)^{2/3}s\big).\] Set \(a_n=2n^{-3}\) and \(T=t=1\). Since \(X(0)>0\) almost surely, \[
 \mathbb P\!\left(\max_{[-1,1]}X_{a_n}>1\right)
 \geq \mathbb P(X(0)>1/n)\longrightarrow1.\] For any fixed positive \(C,c\), the right-hand side of [eq:badFS] (\#eq:badFS) equals \[C2^{2/3}n^{-2}\exp(-c\sqrt2\,n^{-3/2})\longrightarrow0.\] All tested parameters lie within the printed range. This is a contradiction to uniformity in \(a\), rather than a failure of the fixed-tilt Ferrari–Spohn tail estimate.\par
\subsection*{A bounded repair and its limited impact}
Equation reference: eq:repairedFS. The cited fixed-tilt result [Retained latex: \textbackslash{}cite[Lemma\textasciitilde{}5.3, PDF pp.\textasciitilde{}22--23]\{CG\}] requires a time interval of length at least the fixed natural scale. Under scaling, the relevant half-length is \(T'=T(a/2)^{2/3}\), which can be smaller than 1. Apply the fixed-tilt bound on the larger interval of half-length \(\max(1,T')\) and use the scaled threshold \((a/2)^{1/3}t\). After changing absolute constants, this proves the corrected form \[
 \mathbb P\!\left(\max_{[-T,T]}X_a>t\right)
 \leq C(1+Ta^{2/3})\exp(-c\sqrt a\,t^{3/2}).\] Alternatively one can restrict the original form to the range \(T(a/2)^{2/3}\geq1\), or to a uniformly comparable range with constants adjusted.\par
The paper’s coming-down estimates in Corollary 2.8 and Lemma 2.9 impose \(a\geq1\) (PDF pp. 16–17). In the checked large-\(T\) applications, including the heavy-path argument, \(Ta^{2/3}\geq1\), so \(1+Ta^{2/3}\leq2Ta^{2/3}\) and the added constant is absorbed in the prefactor. The condition \(a\geq1\) alone would not justify this absorption for arbitrarily short time intervals; no such unrestricted claim is made here. The classification at arbitrary positive overall tilt also has the separate deterministic scaling reduction. Accordingly, this finding is Minor: the universally quantified auxiliary estimate is false, but no failure of a core conclusion from the small-\(a\) regime is demonstrated. Constants depending on \(a\) would be a different statement and cannot be read into the explicit independence claim.\par
\section*{Coverage, counterevidence and remaining work}
The initial review followed the one-line slope/existence/uniqueness argument in Sections 3–4 and the multiline approximation, confinement and reverse-coupling argument in Sections 5–6. It checked the distinction between a finite interval law and an infinite Gibbs state; endpoint positivity and zero-boundary limiting conventions; the role of asymptotic pinning; tail triviality for extremal states; admissible slope pairs and the two infinite-slope limits; the parabolic normalization; and the order in which \(n\), \(T\), fixed compact windows and small failure probabilities are used.\par
The initial independent algebra file contains 27 passing checks, including four admissible parameter instances, the exact Brownian bridge variance and the small-tilt spatial/time rescaling. These calculations support current Problems 1 and 2 (initial Problems 1 and 3). The withdrawn interval-maximum candidate was not among those checks; its withdrawal rests on direct verification of the source formula, including the factor \(S\).\par
Several apparent local defects were considered but are not additional substantive findings. In particular, the absent additive constant in (5.11) is inconsistent at a zero window, while the stated (5.8) and the relevant reference’s more precise estimate retain the appropriate constant. This is treated as a local transcription issue, not an independent theorem failure. The comparison of an unconditional coupling error with a stronger conditional formulation, and the uniform-in-\(T\) wording near a boundary in (5.5)–(5.6), were also examined. The current report does not establish an additional defect from them. It does not certify every possible reformulation of those steps.\par
For external dependencies, the acquired [CG] PDF has 48 physical pages including its cover; the DOI-bearing article begins on physical page 2. The relevant statements and selected proof passages on mixing, maximum tails, uniform confinement and growing boundary data were read and their evidence pages viewed. For [CIW], the attempted institutional postprint download timed out. A separately acquired 20-page arXiv v1 was used for targeted checks of the tightness, Gibbs-property and monotonicity statements on physical pages 6–8. It is not asserted to be byte-identical or theorem-by-theorem identical to the cited Springer chapter. In particular, its Theorem 1.4 is stated for free and zero boundary ensembles; extending that input to the growing top boundary in the reviewed paper requires the additional compact confinement argument the authors discuss, rather than a verbatim application.\par
Bounded searches of the exact title, correction/erratum terms, the institutional record, the author publication listing and the arXiv record found no applicable correction that resolves these findings. The arXiv record inspected exposes v1 from 10 October 2023; it does not substitute for the frozen institutional source. An author-page fetch timed out, while its indexed listing was available. These observations do not prove that no correction exists. No applicable separate computational supplement was found in the recorded search, and no author code was executed. Missing code is not treated as a mathematical defect.\par
Finding | Established judgment | Repair, adoption and remaining check\par
Problem 1: Major | The printed heavy-path envelope is false, including in its induction regime. | Enlarge the boundary-layer constant and recheck the probability bound and induction. The geometric repair is plausible but not adopted as a completed proof.\par
Problem 2: Minor | The stated all-positive-tilt uniform bound fails as \(a\downarrow0\). | Use [eq:repairedFS] (\#eq:repairedFS) or restrict the scaled interval range. The displayed correction follows from scaling; core applications checked here avoid the failing regime.\par
The report requests correction and verification of the Lemma 6.8 heavy-path estimate. It also records the bounded correction to the small-tilt auxiliary Ferrari–Spohn estimate. The initial interval- maximum finding has been withdrawn: the source states the factor \(S\) in (5.9), and its union-bound derivation is consistent with that form. A false supporting assertion is distinguished from a false final theorem throughout.\par
9 M. Basu Roy Chowdhury, P. Caputo and S. Ganguly, Characterizing Gibbs states for area-tilted Brownian lines, Ann. Probab. 53 (2025), 2196–2255. doi:10.1214/25-AOP1760 (https://doi.org/10.1214/25-AOP1760). Reviewed: the exact frozen institutional 60-page PDF; physical pagination used in this report.\par
P. Caputo and S. Ganguly, Uniqueness, mixing, and optimal tails for Brownian line ensembles with geometric area tilt, Probab. Math. Phys. 6 (2025), 195–239. doi:10.2140/pmp.2025.6.195 (https://doi.org/10.2140/pmp.2025.6.195). Acquired from https://par.nsf.gov/servlets/purl/10610646; targeted pages of the 48-page PDF, including its cover, were used.\par
P. Caputo, D. Ioffe and V. Wachtel, Tightness and line ensembles for Brownian polymers under geometric area tilts, arXiv:1906.06533v1 (15 June 2019). https://arxiv.org/abs/1906.06533v1. The reviewed paper cites the Springer chapter, pp. 241–266, doi:10.1007/978-3-030-29077-1\_10 (https://doi.org/10.1007/978-3-030-29077-1\_10); the acquired version used here is explicitly the arXiv v1.\par
\end{document}
